% ============================================================
% STANDARD CV LAYOUT
% Rendering only; job-specific content lives in ../src/content_CV.tex.
% ============================================================
% Experimental full tagging (testphase=latest) is intentionally not forced here:
% with the current table/ActualText layout it can break compilation or degrade
% ATS text extraction. The PDF still uses PDF management, Unicode mapping and
% explicit ActualText for ATS-critical fields.
% ============================================================
% SHARED COMMANDS
% ============================================================
% Read-only for job-specific content generators.
% Reusable commands shared by all LaTeX projects.

\def\LanguageEnglish{english}
\def\LanguageGerman{ngerman}

% Expandable bilingual selector.
% #1 = English text, #2 = German text.
\newcommand{\MyL}[2]{%
  \ifx\Language\LanguageGerman
    #2%
  \else
    #1%
  \fi
}

% ------------------------------------------------------------
% LLM-READABLE KEY/VALUE STORAGE
% ------------------------------------------------------------
% Keys may contain digits and dots, e.g. Role06.Tasks or Education03.Visible.
% This avoids invalid LaTeX control-sequence names containing digits.
\newcommand{\SetValue}[2]{%
  \expandafter\def\csname SharedValue:#1\endcsname{#2}%
}
\newcommand{\Value}[1]{\csname SharedValue:#1\endcsname}

% Boolean flags use 1 = enabled, 0 = disabled.
\newcommand{\SetFlag}[2]{%
  \expandafter\def\csname SharedFlag:#1\endcsname{#2}%
}
% Expand a flag to its numeric value (1 or 0).
\newcommand{\Flag}[1]{\csname SharedFlag:#1\endcsname}


% Generic conditional for LLM-readable boolean flags.
% #1 = flag key, #2 = enabled branch, #3 = disabled branch.
\newcommand{\IfFlag}[3]{%
  \ifnum\Flag{#1}=1\relax
    #2%
  \else
    #3%
  \fi
}

% ============================================================
% AI-EDITABLE CV CONTENT
% ============================================================
% Source basis:
%   LL_Tabstopp_blaue_Linien(1).dotx
% Job-specific / presentation content only.
% Stable applicant facts: ../src_personal/personal_data.tex
% ============================================================

\def\Language{ngerman}

\ifx\Language\LanguageGerman
  \DocumentMetadata{lang=de-DE,pdfversion=2.0}
\else
  \DocumentMetadata{lang=en-US,pdfversion=2.0}
\fi

% ------------------------------------------------------------
% SECTION VISIBILITY
% ------------------------------------------------------------
\SetFlag{Section.ProfessionalExperience}{1}
\SetFlag{Section.Education}{1}
\SetFlag{Section.Skills}{1}
\SetFlag{Section.Publications}{0}
\SetFlag{Section.Projects}{0}
\SetFlag{Section.ExtraCurricular}{0}
\SetFlag{Section.Interests}{0}

% Three professional stations are present in the Word template.
\SetFlag{Role06.Visible}{0}
\SetFlag{Role05.Visible}{0}
\SetFlag{Role04.Visible}{0}
\SetFlag{Role03.Visible}{1}
\SetFlag{Role02.Visible}{1}
\SetFlag{Role01.Visible}{1}

\SetFlag{Role06.TasksVisible}{0}
\SetFlag{Role05.TasksVisible}{0}
\SetFlag{Role04.TasksVisible}{0}
\SetFlag{Role03.TasksVisible}{1}
\SetFlag{Role02.TasksVisible}{1}
\SetFlag{Role01.TasksVisible}{1}

\SetFlag{Education03.Visible}{1}
\SetFlag{Education02.Visible}{1}
\SetFlag{Education01.Visible}{1}

\SetFlag{Skill.Languages}{1}
\SetFlag{Skill.Set01}{1}
\SetFlag{Skill.Set02}{1}
\SetFlag{Skill.Set03}{0}
\SetFlag{Skill.Set04}{0}
\SetFlag{Skill.Set05}{0}
\SetFlag{Skill.Set06}{0}
\SetFlag{Skill.Set07}{0}
\SetFlag{Skill.Set08}{0}
\SetFlag{Skill.Set09}{0}
\SetFlag{Skill.Set10}{0}

\SetFlag{Publication01.Visible}{0}
\SetFlag{Project01.Visible}{0}
\SetFlag{Project02.Visible}{0}
\SetFlag{Project03.Visible}{0}
\SetFlag{Project04.Visible}{0}
\SetFlag{Activity01.Visible}{0}
\SetFlag{Activity02.Visible}{0}
\SetFlag{Interest01.Visible}{0}
\SetFlag{Interest02.Visible}{0}
\SetFlag{Interest03.Visible}{0}

% ------------------------------------------------------------
% HEADER HIGHLIGHTS
% Terms are taken from the Word CV content.
% ------------------------------------------------------------
\SetValue{Profile.Highlight01}{Bürokaufmann}
\SetValue{Profile.Highlight02}{Büroorganisation}
\SetValue{Profile.Highlight03}{Rechnungsbearbeitung}
\SetValue{Profile.Highlight04}{MS-Office}

% Optional ATS keyword field based only on source wording.
\def\atsKeywords{Bürokaufmann, Büroorganisation, Personalakten, Rechnungsbearbeitung, Kundenbetreuung, Buchführung, MS-Office, Microsoft Word, Microsoft Excel, Microsoft PowerPoint}%

% Education detail fields used by the renderer but not populated by the Word template.
\SetValue{Education03.MainSubjects}{}
\SetValue{Education02.MainSubjects}{}

% ------------------------------------------------------------
% KNOWLEDGE / SKILLS
% ------------------------------------------------------------
\SetValue{Skill.Set01.Content}{%
  \SkillSetBlock{EDV-Kenntnisse}{%
    \CVSkillTabRow{Microsoft Word}{sehr gute Kenntnisse}%
    \CVSkillTabRow{Microsoft Excel}{täglich in Anwendung}%
    \CVSkillTabRow{Microsoft PowerPoint}{Grundkenntnisse}%
  }%
}

\SetValue{Skill.Set02.Content}{%
  \SkillSetBlock{Weiterbildung}{%
    \CVSkillTabRow{04/2017 -- 05/2017: Zeitmanagement-Seminar}{Ausbildungsakademie Muster, München}%
    \CVSkillTabRow{11/2016 -- 12/2016: Kaufmännischer Einsatz für MS-Office}{Ausbildungsakademie Muster, München}%
    \CVSkillTabRow{03/2015 -- 05/2015: Büroorganisation und Statistik}{Ausbildungsakademie Muster, München}%
  }%
}

\SetValue{Skill.Set03.Content}{}%
\SetValue{Skill.Set04.Content}{}%
\SetValue{Skill.Set05.Content}{}%
\SetValue{Skill.Set06.Content}{}%
\SetValue{Skill.Set07.Content}{}%
\SetValue{Skill.Set08.Content}{}%
\SetValue{Skill.Set09.Content}{}%
\SetValue{Skill.Set10.Content}{}%

% ------------------------------------------------------------
% PROFESSIONAL TASKS
% ------------------------------------------------------------
\SetValue{Role03.Tasks}{%
  \begin{itemize}
    \item Erstellung von Präsentationen
    \item Organisation von hauseigenen Messen
    \item Führen und Verwalten von Personalakten
    \item Korrespondenz und Tagesgeschäft
  \end{itemize}%
}

\SetValue{Role02.Tasks}{%
  \begin{itemize}
    \item Eingangsrechnungen kontrollieren
    \item Ausgangsrechnungen erstellen
    \item Angebote unterbreiten und einholen
    \item Forderungsmanagement
  \end{itemize}%
}

\SetValue{Role01.Tasks}{%
  \begin{itemize}
    \item Telefonische Kundenbetreuung
    \item Buchführung und Rechnungsbearbeitung
    \item Steuern und Sozialabgaben abführen
    \item Posteingang
  \end{itemize}%
}

\SetValue{Role04.Tasks}{}%
\SetValue{Role05.Tasks}{}%
\SetValue{Role06.Tasks}{}%

% No project descriptions are present in the Word template.
\SetValue{Project01.DescriptionEN}{}
\SetValue{Project01.DescriptionDE}{}
\SetValue{Project02.DescriptionEN}{}
\SetValue{Project02.DescriptionDE}{}
\SetValue{Project03.DescriptionEN}{}
\SetValue{Project03.DescriptionDE}{}
\SetValue{Project04.DescriptionEN}{}
\SetValue{Project04.DescriptionDE}{}

% ============================================================
% SHARED DOCUMENT FONT-SIZE OPTIONS
% ============================================================
% Loaded before \documentclass by each project entry point.
\def\CVBaseFontSize{12pt}
\def\CoverLetterBaseFontSize{10pt}
\def\ApplicationPackageBaseFontSize{12pt}

\documentclass[\CVBaseFontSize]{article}
\usepackage{graphicx} % Required for inserting images
\usepackage[
    a4paper,
    left=5mm,
    right=5mm,
    top=5mm,
    bottom=5mm,
    bindingoffset=0mm,
    twoside=false
    ]{geometry}
% ============================================================
% PERSONAL SOURCE DATA
% ============================================================
% Source basis:
%   LL_Tabstopp_blaue_Linien(1).dotx
% Contains stable applicant-specific facts and personal asset paths.
% Job-specific generators may READ this file but must not rewrite facts.
% ============================================================

% ------------------------------------------------------------
% IDENTITY AND CONTACT DATA
% ------------------------------------------------------------
\def\ApplicantFirstName{Max}
\def\ApplicantLastName{Mustermann}
\def\ApplicantFullName{\ApplicantFirstName\space\ApplicantLastName}

\def\ApplicantStreet{Musterstraße 12}
\def\ApplicantPostalCode{8080X}
\def\ApplicantCityEN{München}
\def\ApplicantCityDE{München}
\def\ApplicantPhone{089 0123456 / 0177 9912345}
\def\ApplicantEmail{max.mustermann@gmail.com}
\def\ApplicantBirthDate{30. Februar 2001}
\def\ApplicantBirthPlaceEN{Musterstadt}
\def\ApplicantBirthPlaceDE{Musterstadt}
\def\ApplicantMaritalStatus{ledig}

% Personal assets from the CV Word template.
\def\ApplicantPhotoFile{../src_personal/profile_photo.png}
\def\ApplicantSignatureFile{../src_personal/signature.png}

% ------------------------------------------------------------
% PROFESSIONAL ROLE METADATA
% Stable chronology: Role01 = oldest, Role06 = newest.
% Tasks are stored in ../src/content_CV.tex.
% ------------------------------------------------------------

% Role01 | 08.2016 -- 08.2018
\SetValue{Role01.Start}{08/2016}
\SetValue{Role01.End}{08/2018}
\SetValue{Role01.TitleEN}{Sachbearbeiter}
\SetValue{Role01.TitleDE}{Sachbearbeiter}
\SetValue{Role01.AreaEN}{Kundenbetreuung, Buchführung und Rechnungsbearbeitung}
\SetValue{Role01.AreaDE}{Kundenbetreuung, Buchführung und Rechnungsbearbeitung}
\SetValue{Role01.Employer}{Firma GmbH}
\SetValue{Role01.OperationalCompany}{Firma GmbH}
\SetValue{Role01.LocationEN}{München}
\SetValue{Role01.LocationDE}{München}

% Role02 | 09.2018 -- 12.2019
\SetValue{Role02.Start}{09/2018}
\SetValue{Role02.End}{12/2019}
\SetValue{Role02.TitleEN}{Bürokaufmann}
\SetValue{Role02.TitleDE}{Bürokaufmann}
\SetValue{Role02.AreaEN}{Rechnungen, Angebote und Forderungsmanagement}
\SetValue{Role02.AreaDE}{Rechnungen, Angebote und Forderungsmanagement}
\SetValue{Role02.Employer}{Beispiel AG}
\SetValue{Role02.OperationalCompany}{Beispiel AG}
\SetValue{Role02.LocationEN}{München}
\SetValue{Role02.LocationDE}{München}

% Role03 | 01.2020 -- heute
\SetValue{Role03.Start}{01/2020}
\SetValue{Role03.End}{heute}
\SetValue{Role03.TitleEN}{Bürokaufmann}
\SetValue{Role03.TitleDE}{Bürokaufmann}
\SetValue{Role03.AreaEN}{Büroorganisation und Personalakten}
\SetValue{Role03.AreaDE}{Büroorganisation und Personalakten}
\SetValue{Role03.Employer}{Muster GmbH}
\SetValue{Role03.OperationalCompany}{Muster GmbH}
\SetValue{Role03.LocationEN}{München}
\SetValue{Role03.LocationDE}{München}

% Unused role slots intentionally empty.
\SetValue{Role04.Start}{}
\SetValue{Role04.End}{}
\SetValue{Role04.TitleEN}{}
\SetValue{Role04.TitleDE}{}
\SetValue{Role04.AreaEN}{}
\SetValue{Role04.AreaDE}{}
\SetValue{Role04.Employer}{}
\SetValue{Role04.OperationalCompany}{}
\SetValue{Role04.LocationEN}{}
\SetValue{Role04.LocationDE}{}

\SetValue{Role05.Start}{}
\SetValue{Role05.End}{}
\SetValue{Role05.TitleEN}{}
\SetValue{Role05.TitleDE}{}
\SetValue{Role05.AreaEN}{}
\SetValue{Role05.AreaDE}{}
\SetValue{Role05.Employer}{}
\SetValue{Role05.OperationalCompany}{}
\SetValue{Role05.LocationEN}{}
\SetValue{Role05.LocationDE}{}

\SetValue{Role06.Start}{}
\SetValue{Role06.End}{}
\SetValue{Role06.TitleEN}{}
\SetValue{Role06.TitleDE}{}
\SetValue{Role06.AreaEN}{}
\SetValue{Role06.AreaDE}{}
\SetValue{Role06.Employer}{}
\SetValue{Role06.OperationalCompany}{}
\SetValue{Role06.LocationEN}{}
\SetValue{Role06.LocationDE}{}

% ------------------------------------------------------------
% ACADEMIC REFERENCE CONTACTS
% No academic references are present in the Word CV template.
% ------------------------------------------------------------
\SetValue{ReferencePerson01.Name}{}
\SetValue{ReferencePerson01.OrganisationEN}{}
\SetValue{ReferencePerson01.OrganisationDE}{}
\SetValue{ReferencePerson01.DepartmentEN}{}
\SetValue{ReferencePerson01.DepartmentDE}{}
\SetValue{ReferencePerson01.Email}{}
\SetValue{ReferencePerson01.Status}{}
\SetValue{ReferencePerson02.Name}{}
\SetValue{ReferencePerson02.OrganisationEN}{}
\SetValue{ReferencePerson02.OrganisationDE}{}
\SetValue{ReferencePerson02.DepartmentEN}{}
\SetValue{ReferencePerson02.DepartmentDE}{}
\SetValue{ReferencePerson02.Email}{}
\SetValue{ReferencePerson02.Status}{}
\SetValue{ReferencePerson03.Name}{}
\SetValue{ReferencePerson03.OrganisationEN}{}
\SetValue{ReferencePerson03.OrganisationDE}{}
\SetValue{ReferencePerson03.DepartmentEN}{}
\SetValue{ReferencePerson03.DepartmentDE}{}
\SetValue{ReferencePerson03.Email}{}
\SetValue{ReferencePerson03.Status}{}

% ============================================================
% EDUCATION FACTS
% Stable chronology: Education01 = oldest, Education03 = newest.
% ============================================================

% 09.2013 -- 08.2016
\SetValue{Education03.Start}{09/2013}
\SetValue{Education03.End}{08/2016}
\SetValue{Education03.TitleEN}{Abgeschlossene Ausbildung zum Bürokaufmann}
\SetValue{Education03.TitleDE}{Abgeschlossene Ausbildung zum Bürokaufmann}
\SetValue{Education03.GradeEN}{}
\SetValue{Education03.GradeDE}{}
\SetValue{Education03.Institution}{Firma Mustermann GmbH}
\SetValue{Education03.LocationEN}{München}
\SetValue{Education03.LocationDE}{München}
\SetValue{Education03.CourseEN}{}
\SetValue{Education03.CourseDE}{}
\SetValue{Education03.ThesisEN}{}
\SetValue{Education03.ThesisDE}{}
\SetValue{Education03.TermPaperEN}{}
\SetValue{Education03.TermPaperDE}{}

% 09.2012 -- 09.2013
\SetValue{Education02.Start}{09/2012}
\SetValue{Education02.End}{09/2013}
\SetValue{Education02.TitleEN}{Kaufmännisches Berufskolleg}
\SetValue{Education02.TitleDE}{Kaufmännisches Berufskolleg}
\SetValue{Education02.GradeEN}{Abschluss: Fachhochschulreife}
\SetValue{Education02.GradeDE}{Abschluss: Fachhochschulreife}
\SetValue{Education02.Institution}{}
\SetValue{Education02.LocationEN}{}
\SetValue{Education02.LocationDE}{}
\SetValue{Education02.CourseEN}{}
\SetValue{Education02.CourseDE}{}
\SetValue{Education02.ThesisEN}{}
\SetValue{Education02.ThesisDE}{}

% 09.2008 -- 07.2012
\SetValue{Education01.Start}{09/2008}
\SetValue{Education01.End}{07/2012}
\SetValue{Education01.TitleEN}{Mittlere Reife}
\SetValue{Education01.TitleDE}{Mittlere Reife}
\SetValue{Education01.GradeEN}{}
\SetValue{Education01.GradeDE}{}
\SetValue{Education01.Institution}{Realschule Muster}
\SetValue{Education01.LocationEN}{München}
\SetValue{Education01.LocationDE}{München}
\SetValue{Education01.SubjectsEN}{}
\SetValue{Education01.SubjectsDE}{}

% ============================================================
% CONTINUING EDUCATION FROM WORD TEMPLATE
% These source values are rendered through Skill.Set02 in content_CV.tex.
% ============================================================
\SetValue{Training01.Date}{04/2017 -- 05/2017}
\SetValue{Training01.Title}{Zeitmanagement-Seminar}
\SetValue{Training01.Institution}{Ausbildungsakademie Muster}
\SetValue{Training01.Location}{München}

\SetValue{Training02.Date}{11/2016 -- 12/2016}
\SetValue{Training02.Title}{Kaufmännischer Einsatz für MS-Office}
\SetValue{Training02.Institution}{Ausbildungsakademie Muster}
\SetValue{Training02.Location}{München}

\SetValue{Training03.Date}{03/2015 -- 05/2015}
\SetValue{Training03.Title}{Büroorganisation und Statistik}
\SetValue{Training03.Institution}{Ausbildungsakademie Muster}
\SetValue{Training03.Location}{München}

% ============================================================
% PROJECT / PUBLICATION FACTS
% No project/publication content is present in the Word CV template.
% ============================================================
\SetValue{Publication01.StatusEN}{}
\SetValue{Publication01.StatusDE}{}
\SetValue{Publication01.TitleEN}{}
\SetValue{Publication01.TitleDE}{}

\SetValue{Project01.Date}{}
\SetValue{Project01.TypeEN}{}
\SetValue{Project01.TypeDE}{}
\SetValue{Project01.InstitutionEN}{}
\SetValue{Project01.InstitutionDE}{}
\SetValue{Project01.LocationEN}{}
\SetValue{Project01.LocationDE}{}
\SetValue{Project02.Date}{}
\SetValue{Project02.TypeEN}{}
\SetValue{Project02.TypeDE}{}
\SetValue{Project02.InstitutionEN}{}
\SetValue{Project02.InstitutionDE}{}
\SetValue{Project02.LocationEN}{}
\SetValue{Project02.LocationDE}{}
\SetValue{Project03.Date}{}
\SetValue{Project03.TypeEN}{}
\SetValue{Project03.TypeDE}{}
\SetValue{Project03.InstitutionEN}{}
\SetValue{Project03.InstitutionDE}{}
\SetValue{Project03.LocationEN}{}
\SetValue{Project03.LocationDE}{}
\SetValue{Project04.Date}{}
\SetValue{Project04.TypeEN}{}
\SetValue{Project04.TypeDE}{}
\SetValue{Project04.InstitutionEN}{}
\SetValue{Project04.InstitutionDE}{}
\SetValue{Project04.LocationEN}{}
\SetValue{Project04.LocationDE}{}

% ============================================================
% LANGUAGE FACTS
% ============================================================
\SetValue{LanguageSkill01.NameEN}{Deutsch}
\SetValue{LanguageSkill01.NameDE}{Deutsch}
\SetValue{LanguageSkill01.Level}{6}
\SetValue{LanguageSkill02.NameEN}{Englisch}
\SetValue{LanguageSkill02.NameDE}{Englisch}
\SetValue{LanguageSkill02.Level}{4}
\SetValue{LanguageSkill03.NameEN}{Französisch}
\SetValue{LanguageSkill03.NameDE}{Französisch}
\SetValue{LanguageSkill03.Level}{2}

% No extracurricular or interest content is present in the Word CV template.
\SetValue{Activity01.Date}{}
\SetValue{Activity01.RoleEN}{}
\SetValue{Activity01.RoleDE}{}
\SetValue{Activity01.Organisation}{}
\SetValue{Activity01.LocationEN}{}
\SetValue{Activity01.LocationDE}{}
\SetValue{Activity01.DescriptionEN}{}
\SetValue{Activity01.DescriptionDE}{}
\SetValue{Activity02.Date}{}
\SetValue{Activity02.RoleEN}{}
\SetValue{Activity02.RoleDE}{}
\SetValue{Activity02.Organisation}{}
\SetValue{Activity02.LocationEN}{}
\SetValue{Activity02.LocationDE}{}
\SetValue{Activity02.DescriptionEN}{}
\SetValue{Activity02.DescriptionDE}{}

\SetValue{Interest01.TitleEN}{}
\SetValue{Interest01.TitleDE}{}
\SetValue{Interest01.DescriptionEN}{}
\SetValue{Interest01.DescriptionDE}{}
\SetValue{Interest02.TitleEN}{}
\SetValue{Interest02.TitleDE}{}
\SetValue{Interest02.DescriptionEN}{}
\SetValue{Interest02.DescriptionDE}{}
\SetValue{Interest03.TitleEN}{}
\SetValue{Interest03.TitleDE}{}
\SetValue{Interest03.DescriptionEN}{}
\SetValue{Interest03.DescriptionDE}{}

% No additional certificate titles are present in the Word CV template.
\SetValue{AdditionalCertificate01.TitleEN}{}
\SetValue{AdditionalCertificate01.TitleDE}{}
\SetValue{AdditionalCertificate02.TitleEN}{}
\SetValue{AdditionalCertificate02.TitleDE}{}
\SetValue{AdditionalCertificate03.TitleEN}{}
\SetValue{AdditionalCertificate03.TitleDE}{}

% ============================================================
% SHARED VISUAL SYSTEM
% ============================================================
% Read-only for job-specific content generators.
% Central source for font family, colors and reusable font-size commands.

\usepackage[default]{sourcesanspro}
\usepackage{xcolor}

% Core palette
\definecolor{ThemePrimary}{HTML}{0065BD}
\definecolor{ThemeBlack}{HTML}{000000}
\definecolor{ThemeWhite}{HTML}{FFFFFF}
\definecolor{ThemeBlueDark}{HTML}{005293}
\definecolor{ThemeBlueDarker}{HTML}{003359}
\definecolor{ThemeGray80}{HTML}{333333}
\definecolor{ThemeGray50}{HTML}{808080}
\definecolor{ThemeGray20}{HTML}{CCCCCC}
\definecolor{ThemeBeige}{HTML}{DAD7CB}
\definecolor{ThemeOrange}{HTML}{E37222}
\definecolor{ThemeGreen}{HTML}{A2AD00}
\definecolor{ThemeLightBlue}{HTML}{98C6EA}
\definecolor{ThemeMediumBlue}{HTML}{64A0C8}

% Reusable typography
\newcommand{\FontCoverHeaderName}{\fontsize{36pt}{43pt}\selectfont}
\newcommand{\FontCVLastName}{\fontsize{34pt}{30pt}\selectfont}
\newcommand{\FontCVSection}{\large}
\newcommand{\FontPackageSection}{\large\bfseries}

% ============================================================
% AI EDITING BOUNDARY
% ============================================================
% Visibility and job-specific content are read from ../src/content_CV.tex.
% Stable personal and role facts are read from ../src_personal/personal_data.tex.
% This file is rendering-only.
\usepackage{tikz}        % for tikzpicture, \clip, \draw, \node
\usepackage{graphicx}    % for \includegraphics
\usepackage{microtype}
\usepackage{colortbl}
\usepackage{tabularx} % Word-style tab-stop structure for ATS-safe skills
\usepackage{longtable} % page-breakable ATS-safe CV timeline rows
\usepackage{needspace} % keep entry heading + first lines together when possible
\usepackage{array}
\usepackage{enumitem}
\usepackage{placeins} % keep float-based CV entries in source/visual order
\usepackage[main=\Language]{babel}
% \usepackage[main=ngerman, english]{babel}
% \usepackage[main=english, ngerman]{babel}
\usetikzlibrary{calc}
\usepackage[most]{tcolorbox}
\usepackage{ifthen}
\usepackage{pdftexcmds}
% Improve Unicode text extraction for ATS/PDF parsers when using pdfLaTeX.
\input{glyphtounicode}
\pdfgentounicode=1
\usepackage{setspace}
\usepackage{xstring}
\usepackage{etoolbox}
\usepackage[unicode,hidelinks]{hyperref}
\usepackage{accsupp}    % ATS-readable field labels without changing visible banner text

% \usepackage{showframe}
%
% 
% ============================================================
% Variables
% ============================================================
%
% \selectlanguage{\Language}
% --- Document metadata ---
\title{Curriculum Vitae}
\author{\ApplicantFullName}
\date{}
% --- Content ---

% Shared colors and typography are defined in ../helper/style.tex.
\def\arc{1ex}
\def\boxrule{0.2ex}
%
% ============================================================
% ATS-SAFE LAYOUT CONTROL
% ============================================================
% All timeline-like sections use one outer geometry:
%   date/category | decorative vertical rule | content
% The rule is part of each real content row. There is no artificial bottom
% extension, so a block with no tasks ends at its last visible text row.
%
\newlength{\CVSectionBeforeSpace}
\newlength{\CVSectionTitleToRuleSpace}
\newlength{\CVSectionToFirstBlockSpace}
\newlength{\CVBlockBetweenSpace}
\newlength{\CVBulletSize}
\newlength{\CVEntryKeepSpace}
%
\setlength{\CVSectionBeforeSpace}{3.0mm}
\setlength{\CVSectionTitleToRuleSpace}{1.0mm}
\setlength{\CVSectionToFirstBlockSpace}{2.2mm}
\setlength{\CVBlockBetweenSpace}{1.6mm}
% Word template LL_Tabstopp_blaue_Linien.dotx uses a small filled square
% (Wingdings F0A7).  A geometric square is used here so no Wingdings font is
% required and the marker contributes no character to ATS text extraction.
\setlength{\CVBulletSize}{0.55ex}
% Approximate space required for entry title + organization + first task.
\setlength{\CVEntryKeepSpace}{5\baselineskip}
%
\newcommand{\CVAfterBlock}{\par\addvspace{\CVBlockBetweenSpace}}%

% Remove the artificial depth added by LaTeX's table row strut.
% This is always used locally inside a surrounding group/environment.
% The normal strut height is preserved so first-line alignment remains stable,
% while the table/rule no longer extends below the natural content only
% because of \dp\strutbox.
\newcommand{\CVNoRowDepth}{%
  \setbox\strutbox=\hbox{%
    \vrule width 0pt height \ht\strutbox depth 0pt%
  }%
}%
% Current timeline color is locally overwritten by each renderer. This makes
% square task markers automatically use exactly the same color as the rule.
\providecommand{\CVCurrentTimelineColor}{ThemePrimary}%
\newcommand{\CVSquareBullet}{%
  \raisebox{0.16ex}{\textcolor{\CVCurrentTimelineColor}{%
    \rule{\CVBulletSize}{\CVBulletSize}}}%
}%
%
% Bilingual command \MyL is defined in ../helper/common_commands.tex.


% ============================================================
% PERSONAL DATA
% Values are kept in a separate file.
% ============================================================

% Visible PDF text remains unchanged. The ATS/text layer receives a
% descriptive field name immediately before the corresponding value.
\newcommand{\ATSField}[2]{%
    \BeginAccSupp{method=pdfstringdef,unicode,ActualText={#1: #2}}%
    #2%
    \EndAccSupp{}%
}

% Preserve exact punctuation and spacing in the PDF text layer when visible
% formatting changes (e.g. bold labels) would otherwise make extractors insert
% spaces before commas or colons.
\newcommand{\ATSExact}[2]{%
    \BeginAccSupp{method=pdfstringdef,unicode,ActualText={#1}}%
    #2%
    \EndAccSupp{}%
}

% Searchable PDF metadata; values stay synchronized with personal_data.tex.
\hypersetup{
    pdftitle={Curriculum Vitae - \ApplicantFullName},
    pdfauthor={\ApplicantFullName},
    pdfsubject={Curriculum Vitae}
}

\def\meName{%
    {\huge\ATSField{First name}{\ApplicantFirstName}}\\
    {\FontCVLastName\ATSField{Last name}{\ApplicantLastName}}%
}

\def\me{%
    \normalsize\ATSField{Address}{\ApplicantStreet}\\
    \normalsize\ATSField{Postal code and city}{\ApplicantPostalCode\space\MyL{\ApplicantCityEN}{\ApplicantCityDE}}\\
    \normalsize\ATSField{Phone}{\ApplicantPhone}\\
    \normalsize\ATSField{Email}{\ApplicantEmail}\\
    \normalsize\ATSField{Date and place of birth}{\ApplicantBirthDate, \MyL{\ApplicantBirthPlaceEN}{\ApplicantBirthPlaceDE}}%
}

\def\profile{
    \huge \Value{Profile.Highlight01}\\
    \huge \Value{Profile.Highlight02}\\
    \huge \Value{Profile.Highlight03}\\
    \huge \Value{Profile.Highlight04}\\
    }
    
% ============================================================
% COMMAND – Section
% ============================================================
%
\newcommand{\mysection}[2]{%
    % #1 = color
    % #2 = text
    % Section rhythm is controlled only by the global variables above.
    % \addvspace avoids accidental accumulation with the previous block gap.
    \FloatBarrier
    \par\addvspace{\CVSectionBeforeSpace}%
    \noindent
    {\textcolor{#1}{\FontCVSection\MakeUppercase{#2}}}\par
    \nobreak
    \vspace{\CVSectionTitleToRuleSpace}%
    {\color{#1}\hrule height \boxrule}%
    \nobreak
    \vspace{\CVSectionToFirstBlockSpace}%
}%
%
% 
% ============================================================
% COMMAND – Activity
% ============================================================
%
%
% Common parameters
% Spacing between items
% itemsep=2pt → space between items
% parsep=0pt → space within multi-paragraph items
% Space before/after list
% topsep=4pt → space above and below the list
% Indentation / shifting
% leftmargin=1em → moves list left/right
% labelsep=0.5em → space between bullet and text
% labelindent=0pt → extra control of label position
%
\setlist[itemize]{label=\CVSquareBullet,itemsep=1pt,topsep=0pt,leftmargin=1.45em,labelsep=0.55em,partopsep=0pt,parsep=0pt}
%
% Keep the date and the first content line on exactly the same baseline.
% The strut gives both table cells the same reference height even when the
% right-hand text is bold or uses ATS ActualText wrappers.
\newcommand{\CVDateFirstLine}[1]{{\normalsize\strut #1}}
\newcommand{\CVContentFirstLine}[1]{{\strut #1}}
%
\newcommand{\activity}[8]{%
% #1 date start, #2 date end, #3 title, #4 department/grade,
% #5 institution, #6 city, #7 activity content, #8 color
  \begingroup
  \def\CVCurrentTimelineColor{#8}%
  \hyphenpenalty=10000
  \exhyphenpenalty=10000
  \emergencystretch=1.5em
  \Needspace{\CVEntryKeepSpace}%
  \setlength{\LTleft}{0pt}%
  \setlength{\LTright}{0pt}%
  \setlength{\LTpre}{0pt}%
  \setlength{\LTpost}{0pt}%
  \renewcommand{\arraystretch}{1.0}%
  \CVNoRowDepth
  \begin{longtable}{@{}
      p{\CVDateColumnWidth}
      @{\hspace{\CVDateToRuleGap}\color{\CVCurrentTimelineColor}\vrule width \boxrule\hspace{\CVRuleToContentGap}}
      >{\raggedright\arraybackslash}p{\CVContentColumnWidth}
    @{}}
    \CVDateFirstLine{#1 -- #2} &
    \CVContentFirstLine{\ATSExact{#3, #4}{\textbf{#3}, #4}}\par
    {\normalsize #5, #6}\par
    {\normalsize #7}\tabularnewline
  \end{longtable}%
  \CVAfterBlock
  \endgroup
}%

% Professional-experience renderer. The complete task list is rendered
% directly inside the right-hand content cell below the company line.
% No additional task-row or outer bullet wrapper is used.
\newcommand{\professionalactivity}[9]{%
  \begingroup
  \def\CVCurrentTimelineColor{#9}%
  \hyphenpenalty=10000
  \exhyphenpenalty=10000
  \emergencystretch=1.5em
  \IfStrEq{#5}{#6}{%
    \def\EmploymentCompanyLine{%
      {\normalsize
       \ATSExact{\MyL{Employer}{Arbeitgeber}: #5; \MyL{Location}{Ort}: #7}{%
         \textbf{#5}, #7}}%
    }%
  }{%
    \def\EmploymentCompanyLine{%
      {\normalsize
       \ATSExact{\MyL{Contractual employer}{Vertraglicher Arbeitgeber}: #5; \MyL{Operational company}{Einsatzunternehmen}: #6; \MyL{Location}{Ort}: #7}{%
         \textbf{#5}, \MyL{assignment at}{Einsatz bei} #6, #7}}%
    }%
  }%
  \Needspace{\CVEntryKeepSpace}%
  \setlength{\LTleft}{0pt}%
  \setlength{\LTright}{0pt}%
  \setlength{\LTpre}{0pt}%
  \setlength{\LTpost}{0pt}%
  \renewcommand{\arraystretch}{1.0}%
  \CVNoRowDepth
  \begin{longtable}{@{}
      p{\CVDateColumnWidth}
      @{\hspace{\CVDateToRuleGap}\color{\CVCurrentTimelineColor}\vrule width \boxrule\hspace{\CVRuleToContentGap}}
      >{\raggedright\arraybackslash}p{\CVContentColumnWidth}
    @{}}
    \CVDateFirstLine{#1 -- #2} &
    \CVContentFirstLine{\ATSExact{#3, #4}{\textbf{#3}, #4}}\par
    \EmploymentCompanyLine\par
    #8%
    \tabularnewline
  \end{longtable}%
  \CVAfterBlock
  \endgroup
}%
%
% 
% ============================================================
% COMMAND – Skill bars (circle style)
% ============================================================
%
\newcommand{\skill}[3][ThemeBlueDark]{%
  % \noindent
  \begin{tikzpicture}[baseline=(label.base)]
    \pgfmathsetlengthmacro{\R}{0.5ex}
    \pgfmathsetlengthmacro{\D}{2.2ex}

    \node[anchor=west, text width=0.65\linewidth] (label) at (0,0) {#2};
    \coordinate (c) at ($(label.east)+(0.6ex,0)$);

    \foreach \i in {1,...,4}{
      \draw[fill=ThemeGray80!8,draw=ThemeGray80!15]
        ($(c)+(\i*\D,0)$) circle (\R);
    }

    \ifnum#3>0\relax
      \foreach \i in {1,...,#3}{
        \fill[fill=#1]
          ($(c)+(\i*\D,0)$) circle (\R);
      }
    \fi
  \end{tikzpicture}%
}
%
\newcommand{\skillleveltext}[1]{%
  \ifdefstring{\Language}{english}{%
    \ifnum#1=1 Beginner%
    \else\ifnum#1=2 Basic%
    \else\ifnum#1=3 Intermediate%
    \else\ifnum#1=4 Advanced%
    \else\ifnum#1=5 Fluent%
    \else\ifnum#1=6 Native~speaker%
    \else Unknown%
    \fi\fi\fi\fi\fi\fi
  }{%
    \ifdefstring{\Language}{ngerman}{%
      \ifnum#1=1 Anfänger%
      \else\ifnum#1=2 Grundkenntnisse%
      \else\ifnum#1=3 Gute~Kenntnisse%
      \else\ifnum#1=4 Fließend%
      \else\ifnum#1=5 Verhandlungssicher%
      \else\ifnum#1=6 Muttersprache%
      \else Unbekannt%
      \fi\fi\fi\fi\fi\fi
    }{%
      Unknown%
    }%
  }%
}

% #1 English name, #2 German name, #3 integer level
%
% In the skills section, language/tool entries are rendered using the same
% three-tab-stop logic as LL_Tabstopp_blaue_Linien.dotx:
%   category | skill | level
% The content remains one logical top-to-bottom stream for ATS extraction.
\newlength{\CVSkillGroupWidth}
% Left edge of every skill-level label, measured from the beginning of the
% right-hand skill content column. Change only this value to move all levels.
\newlength{\CVSkillLevelStart}
\newlength{\CVSkillCategoryToRuleGap}
\newlength{\CVSkillRuleToSkillGap}
\newlength{\CVSkillCategoryColumnWidth}
\setlength{\CVSkillGroupWidth}{170.10bp}
% Levels are LEFT-aligned at one common position (roughly page centre).
\setlength{\CVSkillLevelStart}{62mm}
% Keep the Word tab stop for the beginning of the skill column. The blue
% separator sits inside the category area, so adding the line does not move
% the skill text to the right.
\setlength{\CVSkillCategoryToRuleGap}{2mm}
\setlength{\CVSkillRuleToSkillGap}{2.5mm}
\setlength{\CVSkillCategoryColumnWidth}{%
  \dimexpr\CVSkillGroupWidth-\CVSkillCategoryToRuleGap-\boxrule-\CVSkillRuleToSkillGap\relax}

\makeatletter
\newcommand{\CVConsumeSkillComma}{%
  \@ifnextchar,{\@gobble}{}%
}
\makeatother

\newcommand{\CVSkillTabRow}[2]{%
  % One logical source row for ATS; only the visual start position of #2 is
  % controlled. Both fields remain ordinary searchable PDF text.
  \par\noindent
  \makebox[\CVSkillLevelStart][l]{#1}%
  \makebox[\dimexpr\linewidth-\CVSkillLevelStart\relax][l]{#2}%
  \par
}

\newcommand{\langskill}[3]{%
  \CVSkillTabRow{\MyL{#1}{#2}}{\skillleveltext{#3}}%
  \CVConsumeSkillComma
}
%
% Skill level text (4 levels)
\newcommand{\skilllevel}[1]{%
  \ifdefstring{\Language}{english}{%
    \ifnum#1=1 Basic%
    \else\ifnum#1=2 Intermediate%
    \else\ifnum#1=3 Advanced%
    \else\ifnum#1=4 Expert%
    \else Unknown%
    \fi\fi\fi\fi
  }{%
    \ifdefstring{\Language}{ngerman}{%
      \ifnum#1=1 Grundkenntnisse%
      \else\ifnum#1=2 Gute~Kenntnisse%
      \else\ifnum#1=3 Fortgeschritten%
      \else\ifnum#1=4 Experte%
      \else Unbekannt%
      \fi\fi\fi\fi
    }{%
      Unknown%
    }%
  }%
}

% #1 English name, #2 German name, #3 integer level
\newcommand{\skillW}[3]{%
  \CVSkillTabRow{\MyL{#1}{#2}}{\skilllevel{#3}}%
  \CVConsumeSkillComma
}
%
% =========================
% Skill Set Renderer
% =========================
%
% #1 group name, #2 skill list
% The first tab stop is the Word-template category column. Rated skills
% generated with \skillW / \langskill use the second and third tab stops.
% Unrated method/format lists use the complete right-hand text area.
\newcommand{\SkillSetBlock}[2]{%
  \begingroup
  \def\CVCurrentTimelineColor{ThemePrimary}%
  \CVNoRowDepth
  \par\noindent
  \begin{tabularx}{\linewidth}{@{}
      p{\CVDateColumnWidth}
      @{\hspace{\CVDateToRuleGap}\color{\CVCurrentTimelineColor}\vrule width \boxrule\hspace{\CVRuleToContentGap}}
      >{\raggedright\arraybackslash}X
    @{}}
    \textbf{#1} & #2\tabularnewline
  \end{tabularx}%
  \CVAfterBlock
  \endgroup
}
%
% === START DOCUMENT ===
%
\begin{document}
%
% ============================================================
% ===== HEADER =====
% ============================================================
%
\newlength{\xPad}
\newlength{\yPad}
\newlength{\ImgH}
\newlength{\BannerH}
\newlength{\AfterBannerH}
\newlength{\PersonalTextGap}
\newlength{\BannerTextEdgeGap}
\newlength{\ApplicantLastNameOpticalShift}
\newlength{\ProfileRightGap}
\newlength{\ProfileBottomEdgeGap}
\newlength{\CVSeparatorPosition}
\newlength{\CVDateToRuleGap}
\newlength{\CVRuleToContentGap}
\newlength{\CVDateColumnWidth}
\newlength{\CVContentColumnWidth}
\newsavebox{\ProfilePhotoBox}
%
\setlength{\xPad}{4mm}
\setlength{\yPad}{3mm}
\setlength{\ImgH}{45mm}
\setlength{\BannerH}{\dimexpr \ImgH + 2\yPad \relax}
\setlength{\AfterBannerH}{\dimexpr \BannerH - 3\yPad + 2mm \relax}
% Horizontal distance between the right edge of the photo and personal data.
% Change only this value to move the whole personal-data block horizontally.
\setlength{\PersonalTextGap}{5mm}
% Shared top/bottom inset for both banner text blocks.
\setlength{\BannerTextEdgeGap}{1.2mm}
% Small optical correction for the large last-name glyph at 34 pt.
\setlength{\ApplicantLastNameOpticalShift}{0mm}
% Right edge of profile highlights mirrors the photo's left page inset.
\setlength{\ProfileRightGap}{\dimexpr \xPad - 0.5mm\relax}
% Optical bottom correction so the last profile line visually reaches the
% same lower boundary as the last personal-data line.
\setlength{\ProfileBottomEdgeGap}{0.5mm}

% Measure the actual photo width at the selected height.  This lets all
% non-Skills vertical separators share the personal-data block's x-position.
\sbox{\ProfilePhotoBox}{\includegraphics[height=\ImgH,keepaspectratio]{\ApplicantPhotoFile}}
% Geometry uses a 5 mm left margin. CVSeparatorPosition is relative to the
% normal text area's left edge; its absolute page x equals personaltext.west.
\setlength{\CVSeparatorPosition}{%
    \dimexpr \xPad + \wd\ProfilePhotoBox + \PersonalTextGap - 5mm\relax}
\setlength{\CVDateToRuleGap}{2mm}
\setlength{\CVRuleToContentGap}{2.5mm}
\setlength{\CVDateColumnWidth}{%
    \dimexpr \CVSeparatorPosition - \CVDateToRuleGap\relax}
\setlength{\CVContentColumnWidth}{%
    \dimexpr \linewidth - \CVSeparatorPosition - \boxrule - \CVRuleToContentGap\relax}
%
\begin{tikzpicture}[remember picture,overlay]
    % Filled top field
    \fill[ThemePrimary]
        (current page.north west)
        rectangle
        ([yshift=-\BannerH]current page.north east);
    
    % Left picture node (rectangular)
    \node[anchor=north west, inner sep=0pt]
        (pic)
        at ([xshift=\xPad,yshift=-\yPad]current page.north west)
        {
        \ifdefempty{\MyBlanc}
            {}
            {\usebox{\ProfilePhotoBox}}
        };

    % Personal data: one logical ATS text block with exactly the same
    % vertical extent as the photo. Each line keeps its natural font height;
    % TeX distributes only the remaining space between the lines.
    \node[
        anchor=north west,
        text=ThemeWhite,
        align=left,
        inner sep=0pt,
        font=\sffamily\bfseries
    ] (personaltext)
        at ([xshift=\PersonalTextGap]pic.north east)
    {
        \ifdefempty{\MyBlanc}
            {}
            {%
            \vbox to \ImgH{%
                \vspace*{\BannerTextEdgeGap}%
                \hbox{{\huge\ATSField{First name}{\ApplicantFirstName}}}%
                \vfill
                \hbox{\hspace*{\ApplicantLastNameOpticalShift}{\FontCVLastName
                    \ATSField{Last name}{\ApplicantLastName}}}%
                \vfill
                \hbox{{\normalsize\ATSField{Address}{\ApplicantStreet}}}%
                \vfill
                \hbox{{\normalsize
                    \ATSField{Postal code and city}{\ApplicantPostalCode\space\MyL{\ApplicantCityEN}{\ApplicantCityDE}}}}%
                \vfill
                \hbox{{\normalsize\ATSField{Phone}{\ApplicantPhone}}}%
                \vfill
                \hbox{{\normalsize\ATSField{Email}{\ApplicantEmail}}}%
                \vfill
                \hbox{{\normalsize
                    \ATSField{Date and place of birth}{\ApplicantBirthDate, \MyL{\ApplicantBirthPlaceEN}{\ApplicantBirthPlaceDE}}}}%
                \vspace*{\BannerTextEdgeGap}%
            }%
            }
    };
    
    % \node[
    %     anchor=north east,
    %     text=ThemeWhite,
    %     inner sep=0pt,
    %     font=\sffamily\bfseries
    % ] (text2)
    %     at ([xshift=-\ProfileRightGap,yshift=-\yPad]current page.north east)
    % {
    %     \vbox to \ImgH{%
    %         \vspace*{\BannerTextEdgeGap}%
    %         \hbox to 0.65\paperwidth{\hfil
    %             {\huge\ATSField{\MyL{Profile highlight 1}{Profil-Highlight 1}}{\Value{Profile.Highlight01}}}}%
    %         \vfill
    %         \hbox to 0.65\paperwidth{\hfil
    %             {\huge\ATSField{\MyL{Profile highlight 2}{Profil-Highlight 2}}{\Value{Profile.Highlight02}}}}%
    %         \vfill
    %         \hbox to 0.65\paperwidth{\hfil
    %             {\huge\ATSField{\MyL{Profile highlight 3}{Profil-Highlight 3}}{\Value{Profile.Highlight03}}}}%
    %         \vfill
    %         \hbox to 0.65\paperwidth{\hfil
    %             {\huge\ATSField{\MyL{Profile highlight 4}{Profil-Highlight 4}}{\Value{Profile.Highlight04}}}}%
    %         \vspace*{\ProfileBottomEdgeGap}%
    %     }%
    % };
\end{tikzpicture}\par
%
%
\vspace*{\AfterBannerH}
%
%
% ============================================================
% Profile
% This section is optional and if you choose to have one it should be customized to each job application. Stick to these guidelines: http://theinterviewguys.com/resume-summary-examples.
% ============================================================
%
% \begin{center}
%     \begin{tcolorbox}[
%         width= 1. \linewidth,
%         colback=ThemePrimary!5,
%         colframe=ThemePrimary,
%         arc=\arc,
%         boxrule=2pt,
%         halign=center,
%         valign=center
%         ]
%         {\textcolor{ThemeBlack}{
% 			{\fontsize{12pt}{14pt}\selectfont{\textbf{\profile}}}
% 			}
%         }
%     \end{tcolorbox}
% \end{center}
%
%
%
%
%
\ifnum\Flag{Section.Education}=1\relax
\mysection{ThemePrimary}{\MyL{Education}{Ausbildung}}

% AI EDUCATION 03 | newest education entry
\ifnum\Flag{Education03.Visible}=1\relax
\activity{\Value{Education03.Start}}{\Value{Education03.End}}
{\MyL{\Value{Education03.TitleEN}}{\Value{Education03.TitleDE}}}
{\MyL{\Value{Education03.GradeEN}}{\Value{Education03.GradeDE}}}
{\Value{Education03.Institution}}{\MyL{\Value{Education03.LocationEN}}{\Value{Education03.LocationDE}}}
{%
\begin{itemize}
  \item \MyL{\Value{Education03.CourseEN}}{\Value{Education03.CourseDE}}
  \item \MyL{Main subject:}{Schwerpunkte:}~\Value{Education03.MainSubjects}
  \item \ATSExact{\MyL{Thesis}{Masterarbeit}:}{\MyL{Thesis}{Masterarbeit}:}~\MyL{\Value{Education03.ThesisEN}}{\Value{Education03.ThesisDE}}
  \item \ATSExact{\MyL{Term paper}{Semesterarbeit}:}{\MyL{Term paper}{Semesterarbeit}:}~\MyL{\Value{Education03.TermPaperEN}}{\Value{Education03.TermPaperDE}}
\end{itemize}%
}{ThemePrimary}%
\fi

% AI EDUCATION 02
\ifnum\Flag{Education02.Visible}=1\relax
\activity{\Value{Education02.Start}}{\Value{Education02.End}}
{\MyL{\Value{Education02.TitleEN}}{\Value{Education02.TitleDE}}}
{\MyL{\Value{Education02.GradeEN}}{\Value{Education02.GradeDE}}}
{\Value{Education02.Institution}}{\MyL{\Value{Education02.LocationEN}}{\Value{Education02.LocationDE}}}
{%
\begin{itemize}
  \item \MyL{\Value{Education02.CourseEN}}{\Value{Education02.CourseDE}}
  \item \MyL{Main subject:}{Schwerpunkte:}~\Value{Education02.MainSubjects}
  \item \ATSExact{\MyL{Thesis}{Bachelorarbeit}:}{\MyL{Thesis}{Bachelorarbeit}:}~\MyL{\Value{Education02.ThesisEN}}{\Value{Education02.ThesisDE}}
\end{itemize}%
}{ThemePrimary}%
\fi

% AI EDUCATION 01 | oldest education entry
\ifnum\Flag{Education01.Visible}=1\relax
\activity{\Value{Education01.Start}}{\Value{Education01.End}}
{\MyL{\Value{Education01.TitleEN}}{\Value{Education01.TitleDE}}}
{\MyL{\Value{Education01.GradeEN}}{\Value{Education01.GradeDE}}}
{\Value{Education01.Institution}}{\MyL{\Value{Education01.LocationEN}}{\Value{Education01.LocationDE}}}
{\MyL{\Value{Education01.SubjectsEN}}{\Value{Education01.SubjectsDE}}}
{ThemePrimary}%
\fi
\fi % Section.Education
\ifnum\Flag{Section.ProfessionalExperience}=1\relax
\mysection{ThemePrimary}{\MyL{Professional Experience}{Berufserfahrung}}

% AI PROFESSIONAL ROLES
% Stable chronology: Role01 = oldest, Role06 = newest.
% Metadata is read from ../src_personal/personal_data.tex.
% Job-specific task lists are read from ../src/content_CV.tex.

% AI ROLE 06 | newest
\ifnum\Flag{Role06.Visible}=1\relax
\professionalactivity{\Value{Role06.Start}}{\Value{Role06.End}}
{\MyL{\Value{Role06.TitleEN}}{\Value{Role06.TitleDE}}}
{\MyL{\Value{Role06.AreaEN}}{\Value{Role06.AreaDE}}}
{\Value{Role06.Employer}}{\Value{Role06.OperationalCompany}}
{\MyL{\Value{Role06.LocationEN}}{\Value{Role06.LocationDE}}}
{\ifnum\Flag{Role06.TasksVisible}=1\relax\Value{Role06.Tasks}\fi}{ThemePrimary}%
\fi

% AI ROLE 05
\ifnum\Flag{Role05.Visible}=1\relax
\professionalactivity{\Value{Role05.Start}}{\Value{Role05.End}}
{\MyL{\Value{Role05.TitleEN}}{\Value{Role05.TitleDE}}}
{\MyL{\Value{Role05.AreaEN}}{\Value{Role05.AreaDE}}}
{\Value{Role05.Employer}}{\Value{Role05.OperationalCompany}}
{\MyL{\Value{Role05.LocationEN}}{\Value{Role05.LocationDE}}}
{\ifnum\Flag{Role05.TasksVisible}=1\relax\Value{Role05.Tasks}\fi}{ThemePrimary}%
\fi

% AI ROLE 04
\ifnum\Flag{Role04.Visible}=1\relax
\professionalactivity{\Value{Role04.Start}}{\Value{Role04.End}}
{\MyL{\Value{Role04.TitleEN}}{\Value{Role04.TitleDE}}}
{\MyL{\Value{Role04.AreaEN}}{\Value{Role04.AreaDE}}}
{\Value{Role04.Employer}}{\Value{Role04.OperationalCompany}}
{\MyL{\Value{Role04.LocationEN}}{\Value{Role04.LocationDE}}}
{\ifnum\Flag{Role04.TasksVisible}=1\relax\Value{Role04.Tasks}\fi}{ThemePrimary}%
\fi

% AI ROLE 03
\ifnum\Flag{Role03.Visible}=1\relax
\professionalactivity{\Value{Role03.Start}}{\Value{Role03.End}}
{\MyL{\Value{Role03.TitleEN}}{\Value{Role03.TitleDE}}}
{\MyL{\Value{Role03.AreaEN}}{\Value{Role03.AreaDE}}}
{\Value{Role03.Employer}}{\Value{Role03.OperationalCompany}}
{\MyL{\Value{Role03.LocationEN}}{\Value{Role03.LocationDE}}}
{\ifnum\Flag{Role03.TasksVisible}=1\relax\Value{Role03.Tasks}\fi}{ThemePrimary}%
\fi

% AI ROLE 02
\ifnum\Flag{Role02.Visible}=1\relax
\professionalactivity{\Value{Role02.Start}}{\Value{Role02.End}}
{\MyL{\Value{Role02.TitleEN}}{\Value{Role02.TitleDE}}}
{\MyL{\Value{Role02.AreaEN}}{\Value{Role02.AreaDE}}}
{\Value{Role02.Employer}}{\Value{Role02.OperationalCompany}}
{\MyL{\Value{Role02.LocationEN}}{\Value{Role02.LocationDE}}}
{\ifnum\Flag{Role02.TasksVisible}=1\relax\Value{Role02.Tasks}\fi}{ThemePrimary}%
\fi

% AI ROLE 01 | oldest
\ifnum\Flag{Role01.Visible}=1\relax
\professionalactivity{\Value{Role01.Start}}{\Value{Role01.End}}
{\MyL{\Value{Role01.TitleEN}}{\Value{Role01.TitleDE}}}
{\MyL{\Value{Role01.AreaEN}}{\Value{Role01.AreaDE}}}
{\Value{Role01.Employer}}{\Value{Role01.OperationalCompany}}
{\MyL{\Value{Role01.LocationEN}}{\Value{Role01.LocationDE}}}
{\ifnum\Flag{Role01.TasksVisible}=1\relax\Value{Role01.Tasks}\fi}{ThemePrimary}%
\fi

\noindent
\fi % CVShowProfessionalExperience

% ============================================================
% PUBLICATIONS
% ============================================================
%
\ifnum\Flag{Section.Publications}=1\relax
\mysection{ThemePrimary}{\MyL{Publications}{Publikationen}}
\begin{table}[h!]
  \CVNoRowDepth
  \centering
  \begin{tabular}{@{}
      p{\CVDateColumnWidth}
      @{\hspace{\CVDateToRuleGap}\color{ThemePrimary}\vrule width \boxrule\hspace{\CVRuleToContentGap}}
      >{\raggedright\arraybackslash}p{\CVContentColumnWidth}
  @{}}
    \ifnum\Flag{Publication01.Visible}=1\relax
      \CVDateFirstLine{\MyL{\Value{Publication01.StatusEN}}{\Value{Publication01.StatusDE}}} &
      \CVContentFirstLine{\MyL{\Value{Publication01.TitleEN}}{\Value{Publication01.TitleDE}}}\\
    \fi
    \vspace{2mm}
  \end{tabular}
\end{table}
\noindent\par
\fi % Section.Publications

% ============================================================
% PROJECTS
% ============================================================
\ifnum\Flag{Section.Projects}=1\relax
\mysection{ThemePrimary}{\MyL{Projects}{Projekte}}
\begingroup
  \hyphenpenalty=10000
  \exhyphenpenalty=10000
  \emergencystretch=1.5em

  \newcommand{\CVProjectEntry}[1]{%
    \begingroup
    \def\CVCurrentTimelineColor{ThemePrimary}%
    \Needspace{4\baselineskip}%
    \setlength{\LTleft}{0pt}%
    \setlength{\LTright}{0pt}%
    \setlength{\LTpre}{0pt}%
    \setlength{\LTpost}{0pt}%
    \CVNoRowDepth
    \begin{longtable}{@{}
      p{\CVDateColumnWidth}
      @{\hspace{\CVDateToRuleGap}\color{\CVCurrentTimelineColor}\vrule width \boxrule\hspace{\CVRuleToContentGap}}
      >{\raggedright\arraybackslash}p{\CVContentColumnWidth}
    @{}}
      \CVDateFirstLine{\Value{#1.Date}} &
      \CVContentFirstLine{\ATSExact{\MyL{\Value{#1.TypeEN}}{\Value{#1.TypeDE}}, \MyL{\Value{#1.InstitutionEN}}{\Value{#1.InstitutionDE}}, \MyL{\Value{#1.LocationEN}}{\Value{#1.LocationDE}}}{%
        \textbf{\MyL{\Value{#1.TypeEN}}{\Value{#1.TypeDE}}}, \MyL{\Value{#1.InstitutionEN}}{\Value{#1.InstitutionDE}}, \MyL{\Value{#1.LocationEN}}{\Value{#1.LocationDE}}}}\par
      \MyL{\Value{#1.DescriptionEN}}{\Value{#1.DescriptionDE}}\tabularnewline
    \end{longtable}%
    \CVAfterBlock
    \endgroup
  }

  \ifnum\Flag{Project01.Visible}=1\relax\CVProjectEntry{Project01}\fi
  \ifnum\Flag{Project02.Visible}=1\relax\CVProjectEntry{Project02}\fi
  \ifnum\Flag{Project03.Visible}=1\relax\CVProjectEntry{Project03}\fi
  \ifnum\Flag{Project04.Visible}=1\relax\CVProjectEntry{Project04}\fi
\endgroup
\noindent\par
\fi % Section.Projects

% ============================================================
% SKILLS
% ============================================================
%
% ATS-safe Word tab-stop layout based on LL_Tabstopp_blaue_Linien.dotx:
%   category | skill | level
% No parallel left/right skill columns. Each category has one blue vertical
% separator before its skill list; the separator uses the same ThemePrimary and
% rule thickness as the earlier CV separators. Source order remains identical
% to the intended ATS reading order.
%
\ifnum\Flag{Section.Skills}=1\relax
\mysection{ThemePrimary}{\MyL{Skills}{Kenntnisse}}
\begingroup
    \setlength{\tabcolsep}{0pt}%
    \renewcommand{\arraystretch}{1.0}%
    \raggedright
    \hyphenpenalty=10000
    \exhyphenpenalty=10000
    \emergencystretch=2em

    \ifnum\Flag{Skill.Languages}=1\relax
    \SkillSetBlock{\MyL{Languages}{Sprachen}}{%
      \langskill{\Value{LanguageSkill01.NameEN}}{\Value{LanguageSkill01.NameDE}}{\Value{LanguageSkill01.Level}},
      \langskill{\Value{LanguageSkill02.NameEN}}{\Value{LanguageSkill02.NameDE}}{\Value{LanguageSkill02.Level}},
      \langskill{\Value{LanguageSkill03.NameEN}}{\Value{LanguageSkill03.NameDE}}{\Value{LanguageSkill03.Level}}%
    }%
    \fi

    \ifnum\Flag{Skill.Set01}=1\relax\Value{Skill.Set01.Content}\fi
    \ifnum\Flag{Skill.Set02}=1\relax\Value{Skill.Set02.Content}\fi
    \ifnum\Flag{Skill.Set03}=1\relax\Value{Skill.Set03.Content}\fi
    \ifnum\Flag{Skill.Set04}=1\relax\Value{Skill.Set04.Content}\fi
    \ifnum\Flag{Skill.Set05}=1\relax\Value{Skill.Set05.Content}\fi
    \ifnum\Flag{Skill.Set06}=1\relax\Value{Skill.Set06.Content}\fi
    \ifnum\Flag{Skill.Set07}=1\relax\Value{Skill.Set07.Content}\fi
    \ifnum\Flag{Skill.Set08}=1\relax\Value{Skill.Set08.Content}\fi
    \ifnum\Flag{Skill.Set09}=1\relax\Value{Skill.Set09.Content}\fi
    \ifnum\Flag{Skill.Set10}=1\relax\Value{Skill.Set10.Content}\fi
\endgroup
\par
%
%
\fi % CVShowSkills
%
% ============================================================
% EXTRA-CURRICULAR
% ============================================================
%
\ifnum\Flag{Section.ExtraCurricular}=1\relax
\mysection{ThemePrimary}{\MyL{Extra-Curricular}{Ehrenamtliches Engagement}}
\begin{table}[h!]
  \CVNoRowDepth
  \centering
  \begin{tabular}{@{}
    p{\CVDateColumnWidth}
    @{\hspace{\CVDateToRuleGap}\color{ThemePrimary}\vrule width \boxrule\hspace{\CVRuleToContentGap}}
    >{\raggedright\arraybackslash}p{\CVContentColumnWidth}
  @{}}
    \ifnum\Flag{Activity01.Visible}=1\relax
      \CVDateFirstLine{\Value{Activity01.Date}} &
      \CVContentFirstLine{\ATSExact{\MyL{\Value{Activity01.RoleEN}}{\Value{Activity01.RoleDE}}, \Value{Activity01.Organisation}, \MyL{\Value{Activity01.LocationEN}}{\Value{Activity01.LocationDE}}: \MyL{\Value{Activity01.DescriptionEN}}{\Value{Activity01.DescriptionDE}}}{%
        \textbf{\MyL{\Value{Activity01.RoleEN}}{\Value{Activity01.RoleDE}}}, \Value{Activity01.Organisation}, \MyL{\Value{Activity01.LocationEN}}{\Value{Activity01.LocationDE}}: \MyL{\Value{Activity01.DescriptionEN}}{\Value{Activity01.DescriptionDE}}}}\\
    \fi
    \ifnum\Flag{Activity02.Visible}=1\relax
      \CVDateFirstLine{\Value{Activity02.Date}} &
      \CVContentFirstLine{\ATSExact{\MyL{\Value{Activity02.RoleEN}}{\Value{Activity02.RoleDE}}, \Value{Activity02.Organisation}, \MyL{\Value{Activity02.LocationEN}}{\Value{Activity02.LocationDE}}: \MyL{\Value{Activity02.DescriptionEN}}{\Value{Activity02.DescriptionDE}}}{%
        \textbf{\MyL{\Value{Activity02.RoleEN}}{\Value{Activity02.RoleDE}}}, \Value{Activity02.Organisation}, \MyL{\Value{Activity02.LocationEN}}{\Value{Activity02.LocationDE}}: \MyL{\Value{Activity02.DescriptionEN}}{\Value{Activity02.DescriptionDE}}}}\\
    \fi
    \vspace{2mm}
  \end{tabular}
\end{table}
\noindent\par
\fi % Section.ExtraCurricular

% ============================================================
% INTERESTS
% ============================================================
%
\ifnum\Flag{Section.Interests}=1\relax
\mysection{ThemePrimary}{\MyL{Interests}{Interessen}}
\begin{table}[h!]
  \CVNoRowDepth
  \centering
  \begin{tabular}{@{}
    p{\CVDateColumnWidth}
    @{\hspace{\CVDateToRuleGap}\color{ThemePrimary}\vrule width \boxrule\hspace{\CVRuleToContentGap}}
    >{\raggedright\arraybackslash}p{\CVContentColumnWidth}
  @{}}
    \ifnum\Flag{Interest01.Visible}=1\relax
      & \textbf{\MyL{\Value{Interest01.TitleEN}}{\Value{Interest01.TitleDE}}}\par
      \MyL{\Value{Interest01.DescriptionEN}}{\Value{Interest01.DescriptionDE}}\\
    \fi
    \ifnum\Flag{Interest02.Visible}=1\relax
      & \textbf{\MyL{\Value{Interest02.TitleEN}}{\Value{Interest02.TitleDE}}}\par
      \MyL{\Value{Interest02.DescriptionEN}}{\Value{Interest02.DescriptionDE}}\\
    \fi
    \ifnum\Flag{Interest03.Visible}=1\relax
      & \textbf{\MyL{\Value{Interest03.TitleEN}}{\Value{Interest03.TitleDE}}}\par
      \MyL{\Value{Interest03.DescriptionEN}}{\Value{Interest03.DescriptionDE}}
    \fi
    \vspace{2mm}
  \end{tabular}
\end{table}
\noindent\par
\fi % Section.Interests

% ============================================================
% SIGNATURE
% ============================================================
% Signature is intentionally outside all section booleans.
%
\begin{tikzpicture}[remember picture,overlay]
    \node[
        anchor=south east,
        xshift=1cm,
        yshift=1.5cm
        ] at (current page.south east) {%
        \begin{tabular}{p{0.55\textwidth} p{0.45\textwidth}}
            & \includegraphics[height=1.8cm]{\ApplicantSignatureFile} \\[-10mm]
            & \rule{7cm}{0.4pt} \\
            & \ATSExact{\MyL{\ApplicantCityEN}{\ApplicantCityDE}, \today}{\MyL{\ApplicantCityEN}{\ApplicantCityDE}, \today}
        \end{tabular}
        };
\end{tikzpicture}
%
%
%
% === END DOCUMENT ===
%
% {\fontsize{1pt}{1pt}\selectfont\color{white}\atsKeywords}%
%
\end{document}